% =============================================================
\section{Discussion}
\label{sec:discussion}
% =============================================================

\subsection{Why does GraphSAGE outperform ResNet at longer horizons?}

At the 6-hour horizon for El~Paso, GraphSAGE-LSTM achieves a 2.2-percentage-point
skill improvement over ResNet-LSTM ($0.636$ vs.\ $0.614$). We hypothesise two
mechanisms:

\textbf{Irregular spatial structure.} At longer lead times, the relevant cloud
signal migrates from the centre of the patch (where it directly causes GHI
reduction) toward the periphery, where incoming cloud advection is reflected in
the satellite channels. CNN filters, being translation-equivariant, apply the same
spatial template across all positions; GraphSAGE, by contrast, learns separate
aggregation rules that can weight distant pixels differently depending on their
graph neighbourhood. This flexibility may be especially valuable for detecting
cloud-edge gradients that precede GHI events.

\textbf{Reduced redundancy.} The ResNet applies multiple stacked convolutions
that effectively redundantly process the spatial information across layers; the
GraphSAGE message-passing is applied once per hop and is less prone to
over-smoothing at the patch scale ($P = 16$). The lower seed variance of
GraphSAGE ($\pm 0.013$ vs.\ $\pm 0.023$ at El~Paso 6h) is consistent with a
more stable loss landscape.

\subsection{Why is forecasting skill lower at Uniandes than at El~Paso?}

Across all deep-learning models and horizons, absolute skill is markedly lower
at the Uniandes mountain site than at El~Paso (e.g., GraphSAGE-LSTM at 6h:
$\skillday = 0.417$ vs.\ $0.636$). We attribute this gap to two related
properties of the Bogot{\'a} climate:

\textbf{Diurnal regularity narrows the room for improvement.}
At $4.6^{\circ}\text{N}$, the solar declination varies little across the year
($\pm 4.6^{\circ}$), producing a highly repeatable clear-sky envelope. Cloud
formation in Bogot{\'a} is predominantly triggered by orographic lifting in
the early afternoon, generating a strong, regular diurnal cycle. Much of the
predictable structure in the Uniandes GHI series is therefore already encoded
in its own recent history and time-of-day, leaving comparatively less
\emph{additional} information for a 4-hour satellite patch sequence to
contribute --- which compresses the achievable skill ceiling for every model,
spatially-aware or not.

\textbf{Satellite signal saturation at longer horizons.} At 6 hours, the
satellite patch observed at $t$ may carry little predictive information about
the ground GHI at $t + 6\,\text{h}$: the cloud system that will eventually
affect the site has often not yet entered the $16\times16$-pixel field of
view, particularly at a site where convection forms locally through
orographic lifting rather than advecting in from a stable upstream direction.
This limits how much spatial encoding --- regardless of architecture --- can
contribute at the longest horizon.

This result suggests that \emph{site-specific climatological structure should
inform both model design and expectation-setting}: in climates with strong,
regular diurnal cloud patterns and locally-forming convection, the marginal
benefit of spatial deep learning over simpler temporal signals is smaller, and
reported skill scores should always be interpreted relative to the local
persistence baseline rather than compared at face value across sites.

\subsection{Why is the 1-hour horizon intractable at El~Paso?}

The negative skill at 1-hour El~Paso (all DL models, $\skillday \approx -0.17$
to $-0.20$) can be explained by the \emph{low persistence error}: at 204.4~W/m²,
persistence is already very accurate because the predominantly clear-sky regime
means irradiance rarely drops by more than a small fraction over 10 minutes.
The DL models, trained to minimise MSE over the full distribution (including
the occasional cloud event), learn a slightly smoothed version of the target
trajectory that incurs higher error on the flat clear-sky majority.

A potential remedy is a loss function that differentially weights cloud events
relative to the clear-sky envelope (e.g., scaling the per-sample MSE by the
departure of observed GHI from the clear-sky reference), or dedicated
cloud-event detection as a preprocessing step.

\subsection{Practical implications}

\begin{itemize}
  \item For \textbf{El~Paso-type sites} (high irradiance, moderate variability),
        GraphSAGE-LSTM with Optuna-tuned hyperparameters is the recommended model
        at 3h and 6h horizons, with a 60+\% reduction in RMSE over persistence.
        At 1h, the system should fall back to persistence or a light-weight
        nowcasting model.

  \item For \textbf{Uniandes-type sites} (tropical mountain, high convective
        variability), even the best models achieve markedly lower skill than at
        El~Paso ($\skillday \leq 0.42$ vs.\ $0.64$ at 6h); operators should
        calibrate reserve margins to the local persistence baseline rather than
        assume that skill levels transfer across climates. GraphSAGE-LSTM remains
        the most consistent top performer at 3h and 6h, but gains over
        persistence are smaller and more sensitive to careful hyperparameter
        tuning.

  \item The high seed variance of ResNet-LSTM at El~Paso 3h ($\pm 0.090$) suggests
        that multiple retraining runs should be averaged for production deployment,
        or that GraphSAGE-LSTM should be preferred for its higher stability.
\end{itemize}

\subsection{Disentangling uncertainty: why aleatoric and epistemic matter separately}
\label{ssec:uq_discussion}

Grid operators require probabilistic outputs --- prediction intervals that
are well-calibrated across conditions of varying cloud cover --- not just
point estimates. Beyond calibration, however, the central premise of this
work is that knowing the \emph{source} of forecast uncertainty is itself
operationally valuable.

\textbf{Why a single symmetric interval is not enough.} A single,
undifferentiated prediction interval --- however it is calibrated --- reports
one half-width per forecast and does not distinguish \emph{why} a given
forecast is uncertain. Two forecasts can carry the same
nominal interval width for entirely different operational reasons: one
because the site is entering a genuinely volatile but well-characterised
convective afternoon (aleatoric), the other because the model has rarely
seen the current satellite-imagery pattern during training (epistemic). An
operator cannot tell these apart from the interval alone, yet the
appropriate response differs --- accept the irreducible variance in the
first case; flag the forecast as low-confidence and prefer a conservative
fallback (e.g.\ persistence, or a wider manually-set margin) in the second.

\textbf{Aleatoric uncertainty as a proxy for cloud-regime volatility.} We
hypothesise that the heteroscedastic variance head (Eq.~\ref{eq:nll}) will
assign larger $\sigma_\text{ale}$ to inputs preceding rapid convective
transitions --- the regime that, on general grounds, is hardest to forecast
deterministically at both sites --- and smaller $\sigma_\text{ale}$ during
stable clear-sky mornings, providing a continuously-varying, input-dependent
signal of intrinsic forecast difficulty that a constant $\hat q$ cannot
express. This is a design hypothesis, not yet an empirical finding
(Section~\ref{sec:discussion}, Limitations); Section~\ref{sec:results} does
not currently attribute forecast errors to specific cloud regimes.

\textbf{Epistemic uncertainty as a data-quality and deployment-readiness
signal.} Because $\sigma_\text{epi}$ (Eq.~\ref{eq:decomposition}) reflects
disagreement among SGLD posterior samples, it is expected to be larger for
satellite-imagery patterns under-represented in the training period (e.g.\
rare cloud morphologies, seasons with sparse data) and to shrink as the
training set grows --- without requiring the variance head itself to change.
This makes $\sigma_\text{epi}$ a natural diagnostic for \emph{when} a
forecasting pipeline should be retrained or supplemented with additional
data at a given site, complementing the purely architecture-level comparison
of Sections~\ref{ssec:resnet}--\ref{ssec:mlp} and~\ref{ssec:ablation}.

\subsection{Implications for solar-integrated power system operation}

The results presented here connect directly to operational decision-making in
solar-integrated power systems.

\textbf{Situational awareness under cloud-induced uncertainty.}
GraphSAGE-LSTM, by learning spatially-aware representations of incoming cloud
systems from satellite imagery, provides operators with earlier detection of
impending irradiance ramps --- the primary source of generation-side
uncertainty in PV-heavy grids. The 63.6\% RMSE reduction over persistence at
6h for El~Paso translates to significantly tighter dispatch intervals and
reduced reserve requirements.

\textbf{Disentangled intervals for risk-stratified dispatch.}
The aleatoric/epistemic decomposition (Section~\ref{ssec:uq}) converts any
point forecast into both an interval and a diagnosis of \emph{why} it is
wide. For a system operator, this supports two distinct, complementary
actions tied to the same forecast: (i) sizing reserve margins to the
combined predictive interval $\bar\mu \pm z_{1-\alpha/2}\sigma_\text{total}$
for routine risk management; and (ii) flagging forecasts with a high
$\sigma_\text{epi}/\sigma_\text{total}$ ratio as candidates for a more
conservative fallback strategy or for prioritised inclusion in future
retraining data, a distinction unavailable from a single undifferentiated
interval.

\textbf{Climate- and horizon-adaptive model selection.}
The marked difference in absolute skill between El~Paso (up to $\skillday =
0.636$ at 6h) and Uniandes (up to $0.417$) has a direct operational
implication: forecasting pipelines should be benchmarked and calibrated
locally rather than assuming that skill levels transfer across climates.
Within a site, the best architecture also depends on the horizon ---
persistence is hard to beat at 1h in the clear-sky-dominated El~Paso regime,
while GraphSAGE-LSTM is the most consistent top performer at 3h and 6h at
both sites. A horizon-aware portfolio strategy --- a light-weight nowcasting
fallback at very short horizons, spatial deep learning at multi-hour horizons
--- could optimise the accuracy-cost tradeoff for operational deployment.

\subsection{Limitations}

\textbf{Geographic scope.} Both sites are located in Colombia and represent
two tropical climate sub-types. Generalisation to other tropical or subtropical
regions requires additional validation.

\textbf{Flat MLP ablation in progress.} The FlatMLP experiments are currently
completing their Optuna runs (16 of the 24 site $\times$ horizon $\times$ seed
configurations evaluated at the time of writing). Final results will be
incorporated to quantify what fraction of deep-learning skill arises from
spatial encoding versus temporal dynamics alone.

\textbf{Graph-construction ablation uses fewer seeds.} The weighted $k$-NN
GraphSAGE-LSTM variant (Table~\ref{tab:ablation}) is evaluated with two
random seeds rather than four, for the compute-budget reasons described in
Section~\ref{ssec:optuna}. This narrower seed budget should be kept in mind
when comparing its mean skill against the four-seed fixed-graph baseline;
the comparison should be treated as indicative pending a larger seed budget.

\textbf{Uncertainty decomposition has not yet been empirically validated.}
The SGLD sampling and variance-head training pipeline (Section~\ref{ssec:uq})
is implemented but the resulting aleatoric/epistemic decomposition
(Table~\ref{tab:decomposition}) has not yet been evaluated on the test set.
In particular, the claim that $\sigma_\text{epi}$ tracks data scarcity and
that $\sigma_\text{ale}$ tracks cloud-regime volatility (Section~\ref{ssec:uq_discussion})
is currently a hypothesis motivating the design, not yet an empirical
finding; both will be assessed quantitatively once results are available.

\textbf{Single-step output.} Current models produce a single-step forecast at
horizon $H$. Simultaneous multi-step output (forecasts for all three horizons
from a single forward pass) could improve consistency and reduce inference cost
in real-time operational forecasting pipelines.
